\documentclass[11pt,a4paper]{article}

\errorcontextlines=9999

\usepackage[T1]{fontenc}
\usepackage[margin=2cm]{geometry}

\usepackage{libertinus}
\usepackage[scale=.95]{inconsolata}

\usepackage[prefix=]{xcolor-material}
\usepackage[colorlinks=true,allcolors=purple]{hyperref}

\usepackage{enumitem}
\usepackage{array}
\usepackage{booktabs}
\usepackage{longtable}
\usepackage{microtype}
\usepackage{parskip}
\usepackage[style={plain number}]{xlistings}
\usepackage{absint}

\LoadLanguages{latex}
\setlist[description]{style=multiline,leftmargin=3cm,font=\bfseries\ttfamily,nosep}
\setlength{\LTpre}{\medskipamount}
\setlength{\LTpost}{\medskipamount}

%% Highlight everything absint defines, taken from the package itself.
\makeatletter
\newcommand*\absintkeywords{%
  \def\absint@kw{}%
  \def\do##1{\appto\absint@kw{\\##1,}}%
  \absintforallsymbols\do
  \absintforallcommands\do
  \expandafter\forcsvlist\expandafter\do\expandafter{\absintapi}%
  \appto\absint@kw{\\absintlegend*,\\mathsf,\\llbracket,\\rrbracket,\\sharp,}%
  % `*' counts as a letter here, the rest are macros this manual mentions
  \edef\@tmp{\noexpand\xlstlangoverride{latex}%
    {morekeywords=[5]{\unexpanded\expandafter{\absint@kw}}}}\@tmp}
\makeatother
\absintkeywords

\newcommand*\cs[1]{\blatex{\\#1}}
\newcommand*\opt[1]{\blatex{#1}}

%% The notation tables are generated by walking the package, see \absintlegend
%% and \absintcommands.
\renewcommand*\absintname[1]{\blatex{\\#1}}
\newcommand*\absinttable[5]{%
  \renewcommand*\absintlegendstart{%
    \begin{longtable}{@{}l@{\quad}l@{\quad}>{\raggedright\arraybackslash}p{#1}@{\quad}>{\raggedright\arraybackslash}p{#2}@{}}
      \caption{#4}\\\toprule \bfseries#3\\\midrule\endfirsthead
      \toprule \bfseries#3\\\midrule\endhead \bottomrule\endfoot}%
  \renewcommand*\absintlegendstop{\end{longtable}}%
  \renewcommand*\absintlegendentry[4]{##1 & ##2 & ##3 & ##4\\}%
  \let\absintcommandentry\absintlegendentry
  #5}

\title{The \texttt{absint} package}
\author{Florian Sihler}
\date{Version 2.0}

\begin{document}
\maketitle

Notation for abstract interpretation: lattices, domains, Galois connections,
widening and narrowing, denotational semantics.

\section{Usage}

\cs{absexpr} switches to absint notation and works in text and in math mode:

\begin{quote}
  \blatex{\\absexpr\{(X, \\partof, \\lub, \\glb, \\bot, \\top)\}}\\
  \absexpr{(X, \partof, \lub, \glb, \bot, \top)}
\end{quote}

For longer parts and for display math use the \texttt{absint} environment:

\begin{minted}{latex}
\begin{absint}
  \[ \Lub_{i \in \N} f^i(\bot) \partof \lfp_{\bot} f \]
\end{absint}
\end{minted}
\begin{absint}
  \[ \Lub_{i \in \N} f^i(\bot) \partof \lfp_{\bot} f \]
\end{absint}

Outside of the two, \cs{lor} and \cs{sqcup} keep their usual meaning.  Names
that are still free (\cs{lub}, \cs{Environments}, \cs{Interval}, \dots) are
defined for plain math mode too, see \opt{stubs}.  \blatex{\\absintuse\{lub\}}
typesets a single symbol by name: \absintuse{lub}.

\section{Concrete and abstract}

\cs{abstract} and \cs{concrete} mark one object:
\absexpr{\abstract{X} \partof \concrete{Y}}.  A whole expression switches
with \cs{abstractly} and \cs{concretely}:

\begin{quote}
  \blatex{\\absexpr\{\\abstractly\{(X, \\partof, \\lub, \\glb)\}\}}\\
  \absexpr{\abstractly{(X, \partof, \lub, \glb)}}
\end{quote}

Only symbols that have an abstract form change; Table~\ref{tab:symbols} shows
both forms of those that do.  \cs{isAbstract} appends the mark to whatever
precedes it, \opt{sharp} chooses it.

\section{Options}\label{sec:options}

Options go to \cs{usepackage}, or later to \cs{absintsetup}:

\begin{minted}{latex}
\absintsetup{color=teal, sharp=\sharp}
\end{minted}

\begin{description}[leftmargin=2.8cm]
  \item[color] colour of the notation.  \opt{purple}.
  \item[colored] colour at all.  \opt{true}.
  \item[mathcolor] colour with \cs{mathcolor} instead of a math group, which
    also catches the scripts that follow a symbol but is markedly slower.
    \opt{false}.
  \item[style] one argument macro applied to every symbol, replacing the
    colouring.  For more, redefine \cs{absintstyle}.
  \item[links] link a symbol to its legend entry.  \opt{true}, off if
    \texttt{hyperref} is missing.  A symbol without an entry stays plain.
  \item[labelprefix] prefix of the placed labels, \opt{absint:} by default.
  \item[track] record the used symbols in the \texttt{.aux} file, needed by
    \cs{absintlegend}.  \opt{true}.
  \item[sharp] mark of the abstract world.  \cs{\#}, often \cs{sharp}.
  \item[flat] mark of the concrete world, appended by \cs{concrete}.  Empty.
  \item[brackets] semantic brackets: \opt{stmaryrd} loads that package,
    \opt{fallback} builds them from two brackets, \opt{auto} takes
    \cs{llbracket} if something else provides it.  \opt{auto}.
  \item[opfont] font of the named operators.  \cs{mathsf}.
  \item[domainfont] font of the domains.  \cs{mathbb}.
  \item[calfont] font of the calligraphic domains.  \cs{mathcal}.
  \item[substitution] how \cs{subst} is set: \opt{assign} gives
    \absexpr{\subst{M}{x}{N}}, \opt{slash} gives
    {\absintsetup{substitution=slash}\absexpr{\subst{M}{x}{N}}}.
    \opt{assign}.
  \item[stubs] define the free names for plain math mode as well.  Also what
    lets \cs{absexpr} survive a \cs{section} title.  \opt{true}.
  \item[shortcut] another name for \cs{absexpr}, without its backslash:
    \opt{shortcut=ai} defines \cs{ai}.
  \item[plain] short for \opt{colored=false, links=false}.
\end{description}

\section{Legend and references}

\cs{absintlegend} lists the notation used in the document and places the
label of each entry, which is what the symbols link to.  It settles after two
to three runs and says when another one is due.  \blatex{\\absintlegend*} lists
everything and \cs{absintcommands} the commands.

\blatex{\\absintmark\{lub,glb\}} places the labels by hand instead; those win
over the legend.  The layout comes from these hooks:

\begin{description}[leftmargin=4.2cm]
  \item[absintlegendstart] before the list.
  \item[absintlegendstop] after the list.
  \item[absintlegendentry] symbol, name, description, aliases.
  \item[absintcommandentry] example, name, description, source.
  \item[absintlegendempty] instead of the list when there is nothing to list.
  \item[absintname] a name.
  \item[absintdemo] the source of an example.
\end{description}

\cs{absintforallsymbols} and \cs{absintforallcommands} apply a one argument
macro to every name.  This manual uses them to highlight the notation and to
lay out Tables~\ref{tab:symbols} and~\ref{tab:commands}.

\section{Lambda calculus}

Abstraction is \cs{lam}, application is \cs{app} and everything else reads
the way it is written:

\begin{minted}{latex}
\absexpr{(\lam{x}{M}) \app N \betared \subst{M}{x}{N}}
\absexpr{\lam{x}{M \app x} \etared M \text{ if } x \notin \FV{M}}
\absexpr{\combY \app f \betareds f \app (\combY \app f)}
\end{minted}

\absexpr{(\lam{x}{M}) \app N \betared \subst{M}{x}{N}}, \quad
\absexpr{\lam{x}{M \app x} \etared M \text{ if } x \notin \FV{M}}, \quad
\absexpr{\combY \app f \betareds f \app (\combY \app f)}

Typing has \cs{typing}, \cs{ctxadd} and the usual arrows, so a rule premise
is one line:

\begin{minted}{latex}
\absexpr{\typing{\ctxadd{\Context}{x}{\sigma}}{M}{\tau}
  \implies \typing{\Context}{\lam{x}{M}}{\sigma \to \tau}}
\end{minted}

\absexpr{\typing{\ctxadd{\Context}{x}{\sigma}}{M}{\tau}
  \implies \typing{\Context}{\lam{x}{M}}{\sigma \to \tau}}

\noindent
and the typed and polymorphic abstractions are \cs{lamty}, \cs{tyabs} and
\cs{forallty}:
\absexpr{\tyabs{X}{\lamty{x}{X}{x}} \hastype \forallty{X}{X \to X}}.

\section{Own notation}

\begin{minted}{latex}
\DeclareAbsintSymbol{signs}{\mathcal{S}}
  [sharp, aliases={Sign}, desc={the sign domain}]
\DeclareAbsintCommand[desc={a widening with a threshold}, demo={\widenat{k}}]
  {widenat}[1]{\widen_{#1}}
\end{minted}

\DeclareAbsintSymbol{signs}{\mathcal{S}}
  [sharp, aliases={Sign}, desc={the sign domain}]
\DeclareAbsintCommand[desc={a widening with a threshold}, demo={\widenat{k}}]
  {widenat}[1]{\widen_{#1}}

\noindent
gives \absexpr{\signs}, \absexpr{\abstractly{\Sign}} and
\absexpr{a \widenat{k} b}.  An empty meaning lifts the command of that name, which
is how the standard symbols are promoted:
\blatex{\\DeclareAbsintSymbol\{lor\}\{\}[bin]}.

\begin{description}[leftmargin=2.2cm]
  \item[class] the math class: \opt{ord}, \opt{bin}, \opt{rel}, \opt{op},
    \opt{open}, \opt{close}, \opt{punct} or \opt{inner}, written as
    \opt{class=bin} or just \opt{bin}.  Keeps the spacing right, so set it.
    \opt{ord}.
  \item[abstract] rendering in the abstract world.
  \item[sharp] short for an \opt{abstract} with the mark appended.
  \item[aliases] further names, sharing rendering and legend entry.
  \item[desc] description for the legend.
\end{description}

\cs{DeclareAbsintAlias} adds a single alias later on.
\cs{DeclareAbsintCommand} is \cs{newcommand}, optional arguments included,
for a command that only exists in absint notation; its \opt{desc} and
\opt{demo} feed \cs{absintcommands}.

\section{The notation}\label{sec:notation}

\absintkeywords% the symbols declared above are part of the notation as well

Symbols with two forms are shown concrete first, abstract second.

\absinttable{.35\textwidth}{.27\textwidth}
  {Symbol & \bfseries Name & \bfseries Description & \bfseries Aliases}
  {The symbols\label{tab:symbols}}{\absintlegend*}

\absinttable{.24\textwidth}{.32\textwidth}
  {Example & \bfseries Name & \bfseries Description & \bfseries Source}
  {The commands\label{tab:commands}}{\absintcommands}

\end{document}
